\subsection{STM32F405}
\label{chap:Back_STM32F405}

In the Embedded Systems I course of the Master's in Electronic Systems at 
the UPM, the STM32F4-DISCOVERY development board is used as part of the 
course materials. This board incorporates an STM32F407VGT6 microcontroller 
with a 32-bit Arm Cortex-M4 processor. However, since the STM32F407VGT6 
is not yet directly implemented in the version of QEMU used in SIVALSE, the 
STM32F405 microcontroller has been selected. It shares many similarities with 
the STM32F407\cite{STM32F405Datasheet2026}\cite{STM32CortexM42026} and, in addition to the I²C protocol and other peripherals, 
already has a solid implementation base. Since this thesis focuses on emulating the 
STM32F405's I²C protocol, this section will list some of its features, taken directly 
from its datasheet.

The STM32F405RG, part of STMicroelectronics' STM32F4 family, is based on 
a 32-bit ARM Cortex-M4 core with a Floating Point Unit (FPU), support for DSP 
instructions, an Adaptive Real-Time Accelerator (ART) for 0-wait execution from 
Flash up to 168 MHz, a Memory Protection Unit (MPU), and 210 DMIPS of 
performance. This configuration makes it an ideal microcontroller for embedded 
systems with real-time and intensive processing requirements, such as motor 
control, power conversion, and multimedia applications. It incorporates up to 1 MB 
of Flash memory, up to 192+4 KB of SRAM (with 64 KB of CCM), and essential 
peripherals: three 12-bit ADCs, two DACs, 17 timers, 16-stream DMA, three I2C ports, 
multiple USART/SPI/CAN ports, Ethernet MAC, USB OTG HS/FS, and up to 
140 fast, 5V-tolerant GPIOs. Its memory architecture maps peripherals between 
0x40000000-0x5FFFFFFFF, as it shows in the table \ref{tab:ARM_memory_map}

\begin{table} [H]
    \centering
    \caption
        {
            STM32F405 Memory Map
        }
    \label{tab:ARM_memory_map}
\end{table}


%The present thesis is centered in the emulation of the STM32F405's I2C protocol. 
%In this section, will be listed some of the STM32F405 features, extracted directly 
%from its datasheet\cite{STM32F405Datasheet2026}\cite{STM32CortexM42026}. This microcontroller is a high-performance 32-bit 
%microcontroller from STMicroelectronics, built around an Arm Cortex-M4 core with FPU, 
%and DSP extensions, and an adaptive real-time accelerator that enables zero-wait-state 
%code execution from Flash at clock frequencies up to 168 MHz. It delivers up 
%to around 210 DMIPS and includes a memory protection unit, making it suitable for 
%complex real-time embedded applications that require both high performance and 
%robust software isolation.
%
%In terms of memory, the STM32F405 integrates up to 1 MByte of embedded Flash 
%and up to 196 KBytes of SRAM, including a 64 KByte core-coupled memory (CCM) 
%region optimized for fast data access, along with OTP memory and a flexible static 
%memory controller for external SRAM, PSRAM, NOR, NAND, and CompactFlash 
%devices. The device operates from 1.8 V to 3.6 V, offers multiple clock sources 
%(external crystal, internal 16 MHz RC, and 32 kHz oscillators for RTC), and supports 
%several low-power modes (Sleep, Stop, Standby) with a dedicated VBAT domain for 
%RTC and backup registers, enabling energy-efficient designs that retain timekeeping 
%and critical data across power cycles.
%
%The STM32F405 features rich analog capabilities, including three 12-bit ADCs 
%with multi-channel and interleaved modes reaching multi-MSPS aggregate sampling 
%rates, as well as two 12-bit DACs for signal generation. It also provides a 
%16-stream general-purpose DMA controller with FIFO and burst support, a CRC unit, 
%a true random number generator, and a large timer subsystem with up to 17 timers 
%(16- and 32-bit) supporting capture/compare, PWM generation, pulse counting, 
%and quadrature encoder interfaces, making it well suited for motor control, 
%power conversion, and advanced control applications.
%
%On the digital I/O and connectivity side, the STM32F405 exposes up to 140 
%GPIOs (most of them fast and 5 V tolerant) and supports a wide range of communication 
%peripherals: up to three I2C controllers, multiple USART/UART interfaces 
%with support for protocols like LIN, IrDA, and ISO 7816, up to three SPI/I2S blocks, 
%two CAN 2.0B controllers, SDIO, and a parallel camera interface. Advanced connectivity 
%is further enhanced by an integrated Ethernet MAC with IEEE 1588v2 
%hardware timestamping and DMA, as well as full-speed and high-speed USB 2.0 
%OTG controllers with on-chip PHY and dedicated DMA, enabling the microcontroller 
%to serve as a communication hub in networked and multimedia embedded 
%systems.
%
%For development and system integration, the MCU includes a serial-wire/JTAG 
%debug interface, an Embedded Trace Macrocell for instruction tracing, an LCD 
%parallel interface, a real-time clock with calendar and subseconds, and a 96-bit 
%unique device ID that can be used for traceability or security-related features.